\section{Explicit level and spin cutoffs for orthogonal normal cores}
\label{sec:orthogonal-level-cutoff}

\paragraph{Status and scope.}
This is a new author derivation from published quantitative character
estimates. The complete proof, invoked sources and author arithmetic certificate
have been read by /root and /root/g34\_eightfolds. Both whole-proof
reviews pass. Separately authored exact arithmetic checks also pass;
their actual scopes and timings are retained in the check register.
This section gives a strict reduction of the orthogonal normal-core
gate. It does not give an upper bound on the natural dimension of
the remaining ordinary low-level orthogonal representations.

Write $d$ for the natural orthogonal dimension, $a$ for the degree
of a complex irreducible representation, and put
\[
 \sigma_0=10^{-29},\qquad J_0=10^{30},\qquad
 M_0=2\cdot10^{29}=J_0/5,\qquad D_0=2\cdot10^{30}.
                                                               \tag{OL.1}
\]
For odd $d$ we assume $q$ odd. For even $d$ either characteristic
is allowed, with the appropriate sign $\epsilon\in\{+,-\}$.
For a finite-order operator $T$ on a complex space of dimension $a$,
let
\[
                 \nu(T)=a-\max_\lambda\dim\ker(T-\lambda).
                                                               \tag{OL.2}
\]
This number is unchanged by scalar multiplication of $T$.

\begin{lemma}[D-OL-01: the quantitative orthogonal inputs]
\label{D-OL-01}
Let $d\ge10$. Let $S=\Omega_d^\epsilon(q)$, and for odd $q$
also allow $\widetilde S=\operatorname{Spin}_d^\epsilon(q)$.
The sign is suppressed when $d$ is odd. Every ordinary
irreducible character $\chi$ of either indicated group satisfies
\[
 \frac{|\chi(g)|}{\chi(1)}
 \le \chi(1)^{-\sigma_0\operatorname{supp}(g)/d}.
                                                               \tag{OL.3}
\]
For a spin element, support is computed on its image in the
natural orthogonal group.

If $\chi\in\operatorname{Irr}(S)$ has level $\ell$, then
\[
 0\le\ell\le\lfloor d/2\rfloor,\qquad
 \chi(1)\le\left(\frac{q^d-1}{q-1}\right)^\ell<q^{d\ell}
 \quad(\ell>0).
                                                               \tag{OL.4}
\]
For $d\ge50$ and $k=\lceil\ell/3\rceil$,
\[
              \chi(1)\ge q^{k(d-2k-2)}(q-1)^k.
                                                               \tag{OL.5}
\]

For odd $q$, the following two failures both imply
\[
                 a>q^{E(d)},\qquad
                 E(d)=d^2/8-d/4-2:
                                                               \tag{OL.6}
\]
an irreducible spin character of degree $a$ does not descend to
$S$; or an irreducible character of $S$ of degree $a$ does not
extend to $\operatorname{SO}_d^\epsilon(q)$.
\end{lemma}

\begin{proof}
Equation (OL.3) is Larsen--Tiep's published uniform character
theorem, with the smaller explicit constant $\sigma_0$ already
checked in D-UNW-01. Its exact source is
\cite[Theorem~5.5 and complete proof, accepted-author pp.~25--26]
{UNW:LT2024}. That theorem explicitly includes $\Omega(V)$ and,
in odd characteristic, $\operatorname{Spin}(V)$. The support
convention for the latter is also stated immediately before
Proposition~6.4, author p.~30. Its larger available constant is
$1/(6507\cdot2^{21}\cdot10^{18})>7\cdot10^{-29}$.

We use exactly the level of
\cite[Definition~3.2(ii), p.~160]{OL:GLT2024}: the least exponent
of $\tau_A+\zeta_A$ containing the character. Here
$\tau_A(g)=q^{\dim\ker(g-1)}$ and
$\zeta_A(g)=(-1)^d(-q)^{\dim\ker(g-1)}$ on the natural
quadratic space $A$. Lemma~3.3 proves that these two characters
agree on $\Omega(A)$ in all the characteristics at issue.
Lemma~3.4(iii), including its proof on pp.~160--161, gives the
range in (OL.4). Theorem~1.5(iii), with its proof at the end of
Section~3 on p.~167, gives the degree upper bound.
Theorem~3.7 and its complete proof, pp.~165--167, give (OL.5):
for $\ell>0$ we have $\ell>3(k-1)$, so its contrapositive applies;
the case $\ell=0$ is immediate. The small
exceptional pairs in that theorem are absent for $d\ge50$.

For the spin assertion, write $d=2m+1$ or $2m$.
Proposition~6.3 and its complete proof, pp.~193--194 of
\cite{OL:GLT2024}, give respectively
\[
 a>q^{m(m+1)/2}/4,\qquad a>q^{m(m-1)/2}/4.
                                                               \tag{OL.7}
\]
In the odd-dimensional case the kernel of the spin map is the
entire center of order two. A character nontrivial on that
kernel is faithful because the spin group is quasisimple, so
the faithfulness hypothesis of Proposition~6.3(i) is satisfied.
In even dimension Proposition~6.3(ii) directly has the required
noninflation hypothesis; faithfulness on the entire center is
not imposed. Since $4\le q^2$, both estimates imply (OL.6).

For the extension assertion, $S$ has index two in
$\operatorname{SO}_d^\epsilon(q)$. An invariant irreducible
character of $S$ extends across this cyclic quotient: an
intertwiner for a coset representative can be rescaled so that
its square is the prescribed image of the square of that
representative. Thus if the character does not extend, its
induction is irreducible and has degree $2a$.
Proposition~7.3 and its complete proof, pp.~198--200 of
\cite{OL:GLT2024}, apply to this induced character and give
\[
 2a>q^{m^2/2}\quad(d=2m+1),\qquad
 2a>(q-1)q^{m(m-1)/2-1}\quad(d=2m).
                                                               \tag{OL.8}
\]
Using $2\le q$ and $q-1\ge1$, these imply (OL.6).
Indeed $m^2/2-1=E(2m+1)+9/8$, while the weaker exponent
$m(m-1)/2-2$ in even dimension is exactly $E(2m)$.
All propositions used here are ordinary-character statements;
no cross-characteristic modular variant is needed.
\end{proof}

\begin{lemma}[D-OL-02: an explicit degree barrier]
\label{D-OL-02}
Suppose a finite group $J\subset\operatorname{GL}(V)$ normalizes
an irreducibly acting central product of quasisimple groups
$E_1\cdots E_t$. On each associated tensor factor, suppose the
representation lifts to an ordinary irreducible character of
$\operatorname{Spin}_{d_i}^{\epsilon_i}(q_i)$ for odd $q_i$, or
$\Omega_{d_i}^{\epsilon_i}(q_i)$ for even $q_i$, with $d_i\ge10$.
Let its degree be $a_i$, and assume for every factor that
\[
                              a_i\ge q_i^{M_0d_i}.
                                                               \tag{OL.9}
\]
Then every nonscalar element $T\in J$ satisfies
\[
                           \nu(T)\ge\frac3{16}\dim V.
                                                               \tag{OL.10}
\]
Such a representation cannot be an inducing block of an actual
finite-cover Weil sixfold source with period rank nine.

In particular, in the transitive normal-core situation of
D-BI-19 with orthogonal simple factors of natural dimension
$d\ge10$ and field size $q$, every factor degree necessarily
satisfies
\[
                                  a<q^{M_0d}.
                                                               \tag{OL.11}
\]
All compatible outer automorphisms, represented central
quotients, scalar extensions and tensor-factor permutations
are included.
\end{lemma}

\begin{proof}
For one factor, a noncentral element of the covering group has
natural support at least one. In the spin case, an element
whose natural image is scalar is central: its commutators lie
in the central kernel of the natural map, and the resulting
homomorphism from the perfect group to that kernel is trivial.
Equations (OL.3) and (OL.9) therefore imply
\[
                         |\chi(g)|\le a q^{-2}\le a/4
                         \quad(g\notin Z).
                                                               \tag{OL.12}
\]

Let $T$ be a nonscalar compatible linear normalizer of this
irreducible factor. By Schur's lemma there is a core operator
$x$ with $b=TxT^{-1}x^{-1}\ne1$. Choose an eigenvalue
$\lambda$ of $T$ with largest eigenspace. Then
\[
 \operatorname{rank}(b-1)
   =\operatorname{rank}(Tx-xT)
   \le2\operatorname{rank}(T-\lambda)=2\nu(T).
                                                               \tag{OL.13}
\]
If $b$ is a nonidentity scalar its moving rank is $a$.
Otherwise lift it to a noncentral element of the covering
group and use (OL.12). Since $b$ belongs to a finite group,
its eigenvalues have absolute value one, and
\[
 2\operatorname{rank}(b-1)
 \ge a-\operatorname{Re}\operatorname{Tr}(b)\ge3a/4.
\]
Together these give $\nu(T)\ge3a/16$. This argument does not
require a classification of the automorphism implemented by
$T$, or that its coefficient field equal the quadratic Weil
field.

An operator fixing all tensor factors is a tensor product of
the individual intertwiners. The maximal eigenspace proportion
of a tensor product is at most that of any selected factor,
by convolution of eigenvalue multiplicities. If the total
operator is nonscalar, some individual intertwiner is
nonscalar and gives (OL.10).

For an actual permutation cycle of length $s\ge2$, the factor
degrees along that cycle coincide, say they equal $a\ge2$.
Changes of tensor bases reduce the operator to a weighted
cyclic shift with one return operator. The return operator is
diagonalizable because a power is scalar, as follows from the
finite projective order of the total operator. In its
eigenbasis each orbit of tensor words contributes at most one
copy of a prescribed eigenvalue. There are $a$ constant words,
and every nonconstant orbit has size at least two. Thus every
eigenspace has dimension at most
\[
                       a+(a^s-a)/2\le3a^s/4.
\]
Consequently such a cycle gives $\nu\ge a^s/4$, and convolution
with the other factors preserves this proportion. Since
$1/4>3/16$, equation (OL.10) follows in all cases.

We recall the actual geometric constraint. For an inducing
degree $e=\dim V$, let $X$ be the intermediate curve of genus
$h$ and let $r$ be the number of nontrivial inertia points.
Polarized Shapiro and the finite-cover factorization already
proved in this source problem give
\[
 6=2e(h-1)+\sum_{j=1}^r\operatorname{rank}(g_j-1).
                                                               \tag{OL.14}
\]
The additional finite-cover data over pointed $X$ are discrete,
so the period map factors through a finite cover of its moduli.
Here $e>24$ by (OL.9). Thus $h\ge2$ is impossible from
(OL.14). If $h=1$, then $r\le6$, and pointed elliptic-curve
moduli has dimension at most six (or one when $r=0$), so cannot
support period rank nine. Hence rank nine requires $h=0$,
$r\ge12$, and total moving rank $2e+6$.
Every nonidentity scalar inertia has moving rank $e$; every
nonscalar inertia has moving rank at least $3e/16$. Therefore
\[
                   2e+6\ge12\cdot3e/16=9e/4,
                   \qquad e\le24,
                                                               \tag{OL.15}
\]
a contradiction. This is an obstruction from the actual
factorization and period rank, not a claimed sufficiency of a
dimension inequality.

Finally, at $d\ge10$ there is no exceptional Schur multiplier
for these orthogonal simple groups; the universal covers are
the indicated spin groups in odd characteristic and the
ordinary orthogonal groups in even characteristic. This is
the precise exceptional-cover list in
\cite[Chapter~5, Table~1, p.~35; ordinary-cover discussion,
p.~4]{BI:HHM}. All quasisimple factors are central quotients
of those covers. Transitive conjugation of the simple factors
makes their degrees equal, and their abstract groups have the
same orthogonal parameters at these dimensions. If one common
degree violated (OL.11), every factor would satisfy (OL.9),
and the preceding obstruction would apply. This proves
(OL.11). In particular, no assertion that character level is
invariant under an outer automorphism is needed for this step.
\end{proof}

\begin{theorem}[D-OL-03: the remaining orthogonal gate]
\label{D-OL-03}
Consider the transitive orthogonal normal-core situation of
D-OL-02 for an actual period-rank-nine Weil sixfold source.
Every factor character which descends to
$\Omega_d^\epsilon(q)$ has
\[
                                  \ell<10^{30}.
                                                               \tag{OL.16}
\]
If $d\ge2\cdot10^{30}$, every factor character descends to
$\Omega_d^\epsilon(q)$. When $q$ is odd, every such descended
character also extends to $\operatorname{SO}_d^\epsilon(q)$.
Equivalently, any genuinely spin factor, or any ordinary factor
failing this special-orthogonal extension, has
\[
                                  d<2\cdot10^{30}.
                                                               \tag{OL.17}
\]
These conclusions require neither a self-duality assumption
nor a prescribed determinant class.
\end{theorem}

\begin{proof}
If $\ell\ge J_0$, then (OL.4) forces $d\ge2J_0$, in
particular $d\ge50$. Put $k=\lceil\ell/3\rceil$.
The inequalities
\[
 \ell/3\le k\le\ell/3+2/3\le d/6+2/3
\]
give
\[
 d-2k-2\ge2d/3-10/3
             =3d/5+(d-50)/15\ge3d/5.
\]
Consequently (OL.5) implies
\[
                      a\ge q^{3dk/5}\ge q^{d\ell/5}
                        \ge q^{dJ_0/5}=q^{M_0d},
                                                               \tag{OL.18}
\]
contradicting (OL.11). This proves (OL.16).

For the remaining assertions write $d=2J_0+s$ with $s\ge0$.
An exact expansion gives
\[
 \begin{split}
 E(d)-M_0d
  &=J_0^2/10-J_0/2-2
      +(3J_0/10-1/4)s+s^2/8>0.
 \end{split}
                                                               \tag{OL.19}
\]
The constant term is positive already for every integer
$J_0\ge10$, and all remaining coefficients are positive.
Either failure in (OL.6) would therefore imply
$a>q^{M_0d}$, again contradicting (OL.11).
For even $q$ the high-dimensional universal covering group
already is the ordinary orthogonal group, so only the descent
assertion is relevant; no odd-characteristic extension
proposition is applied there.
\end{proof}

\paragraph{Precisely what remains.}
In the nonsplit primitive normal-core setting, the previous
field theorem D-BI-20 also gives $q\le27$; that additional
hypothesis is retained when using this field cutoff. The new
result leaves an ordinary orthogonal gate of bounded character
level $\ell<10^{30}$ but unbounded natural dimension $d$.
For $d\ge2\cdot10^{30}$ it consists of characters descending
to $\Omega$ and, in odd characteristic, extending to
$\operatorname{SO}$. This is a strict level/cover reduction,
not a finite classification of all orthogonal sources.
The highest-level theta constituent in
\cite[Corollary~4.7, p.~174]{OL:GLT2024} is proved nonzero;
that statement is not used as an irreducibility or exhaustion
theorem. No extension to the full orthogonal group is asserted.

\paragraph{Exact certificate and source boundary.}
The standalone program \texttt{orthogonal\_levels/check\_constants.py}
checks the rational exponent constants, the polynomial identities
in (OL.18)--(OL.19), the four spin/extension exponent reductions,
and finite boundary instances of the level inequality. It is
an arithmetic certificate for the displayed universal proof,
not an enumeration of enormous finite groups. Complete invoked
published proofs and the accepted Larsen--Tiep proof were read;
no new proof of their classification or Deligne--Lusztig inputs
is claimed. The source log preserves exact versions and hashes.

